\documentclass[tikz,border=10pt]{standalone}
\usepackage{ctex}
\usepackage{amsmath}
\usetikzlibrary{calc, arrows.meta}

\begin{document}
% 线面垂直的向量判定示意图
% 直线 l 垂直平面 α ⟺ 方向向量 a ∥ 法向量 n
\begin{tikzpicture}[
  x={(1.0cm,0cm)},
  y={(0.40cm,0.30cm)},
  z={(0cm,0.92cm)},
  thick,
  dot/.style={fill=black, circle, inner sep=1.5pt},
]

% 画平面 α（平行四边形）
\fill[blue!12, draw=blue!50, fill opacity=0.75]
  (-2.0,0,0) -- (2.0,0,0) -- (2.8,3.0,0) -- (-1.2,3.0,0) -- cycle;
\node[blue!70, font=\small\bfseries] at (2.5, 1.5, 0) {$\alpha$};

% 平面内两不共线向量 u, v（帮助理解法向量概念）
\coordinate (O0) at (0, 1.5, 0);
\coordinate (Pu) at (1.5, 1.5, 0);
\coordinate (Pv) at (0.3, 2.8, 0);
\draw[blue!60, -Stealth] (O0) -- (Pu) node[right, font=\scriptsize, blue!70] {$\vec{u}$};
\draw[blue!60, -Stealth] (O0) -- (Pv) node[above, font=\scriptsize, blue!70] {$\vec{v}$};

% 垂足 A 在平面内
\coordinate (A) at (0, 1.5, 0);
\node[dot] at (A) {};
\node[font=\small, below left] at (A) {$A$};

% 法向量 n（垂直平面，向上）
\draw[orange, very thick, -Stealth] (A) -- (0, 1.5, 1.8)
  node[right, font=\small\bfseries, orange] {$\vec{n}$};

% 直线 l 上一点 B（在平面外）
\coordinate (B) at (0, 1.5, 2.5);
\node[dot] at (B) {};
\node[font=\small, right] at (B) {$B$};

% 直线 l（从平面下方穿过）
\coordinate (Lbot) at (0, 1.5, -0.5);
\draw[red, thick] (Lbot) -- (B);
\node[red, font=\small, right] at ($(Lbot)+(0.1,0)$) {$l$};

% 方向向量 a ∥ n（与 n 平行，用红色标出）
\draw[red, very thick, -Stealth] (A) -- (0, 1.5, 1.4)
  node[left, font=\small\bfseries, red] {$\vec{a}$};

% 平行符号提示
\node[font=\small, align=center] at (1.8, -0.5, 1.2)
  {$\vec{a} \parallel \vec{n}$};

% 垂直符号（在 A 处）
\draw[gray, thin]
  ($(A)+(0.12,0,0)$) -- ($(A)+(0.12,0,0.12)$) -- ($(A)+(0,0,0.12)$);

% 公式框
\node[draw=red!50, rounded corners, fill=red!5, font=\small, align=center,
      inner sep=4pt] at (3.5, 4.5, 0)
  {$l \perp \alpha$\\[2pt]
   $\Leftrightarrow\;\vec{a}\parallel\vec{n}$\\[2pt]
   $\Leftrightarrow\;\vec{a}=\lambda\vec{n}$};

% 底部说明
\node[font=\small, align=center] at (0.5, -0.8, 0)
  {线面垂直：方向向量 $\vec{a}$ 平行于法向量 $\vec{n}$};

\end{tikzpicture}
\end{document}
